\textbf{Задача 32} Определим класс $\bpp/_{\mathrm{poly}}$ как множество языков $A$, для которых существует семейство схем из функциональных элементов $\{C_n\}$ с двумя группами входов $x$ и $r$ со следующим свойством: если $x\in A$, то $\mathrm{Pr}_r[C_n(x,r)=1]>\frac{2}{3}$, а если $x\not\in A$, то $\mathrm{Pr}_r[C_n(x,r)=1]<\frac{1}{3}$. Докажите, что $\bpp/_{\mathrm{poly}}=\p/_{\mathrm{poly}}$.

\begin{proof} Пусть $A\in\bpp/_{\mathrm{poly}}$ произвольный, распознается семейством вероятностных схем $\{C_n\}$.
Зафиксируем $n$. Можно за время $O(p(n))$ сделать амплификацию, и уменьшить ошибки до $<2^{-n}$. Обозначим полученную схему $C_n'$.

Тогда $\forall x\in\{0,1\}^n$ доля $r$, которые ошибаются на этом $x$, меньше $2^{-n}$. Значит, доля $r$, которые ошибаются хотя бы на одном $x$, меньше $1$. Поэтому существует $r$ полиномиальной длины, для которого $C_n'$ не ошибается при любом входе длины $n$. Зашьём его в схему $C_n'$. Получим, полиномиальную схему $C_n''$, распознающую $A$ на входе длины $n$. 

Сделаем так для каждого $n$ и получим, что $A\in\p/_{\mathrm{poly}}$. В силу произвольности $A$, $\bpp/_{\mathrm{poly}}=\p/_{\mathrm{poly}}$. \end{proof} 